Le $\mathrm{pH}$ d'une solution aqueuse d'ibuprofène de concentration molaire
$C=\qty{5.0e-2}{\mol\per\L}$ vaut $\mathrm{pH}=2,7$ à
$\qty{25}{\degreeCelsius}$.

L'équation de la réaction modélisant la transformation entre l'ibuprofène et
l'eau est~:
\[
    \ce{C13H18O_{2(aq)} + H2O_{(\ell)} <=> C13H17O^-_{2(aq)} + H3O^+_{(aq)}}
\]

\begin{enumerate}
    \item Montrer que cette transformation est limitée.
    \item Calculer la valeur du quotient de réaction $Q_{r,\eq}$ du système
          chimique à l'équilibre.
    \item En déduire la valeur du $\mathrm{pK_{A}}$ du couple
          $(\ce{C13H18O_{2(aq)}}\ttslash{}\ce{C13H17O^-_{2(aq)}})$.
\end{enumerate}

